\scanpage{18}
\mspage{27}
\mspage{27}
\deleted{
\uline{Corollaire} 5.9. Les foncteurs 5.7. induisent des anti-équivalences
quasi-inverses l'une de l'autre entre la catégorie des groupes de type
multiplicatif de type fini\inserted{$H$} sur $S$, et la catégorie des groupes constants
tordus à engendrement fini\inserted{$R$} sur $S$. Tout tel groupe $R$ est quasi-isotrivial.

\uline{Corollaire} 5.10. Soient $H,G$ deux \deleted{\unsure{de}} $S$-préschémas en groupes de type
multiplicatif et de type fini. Alors $\underline{\operatorname{Isom}}_{S\text{-gr}}(H,G)$ est représentable
par un sous-préschéma ouvert\inserted{et fermé} de $\underline{\operatorname{Hom}}_{S\text{-gr}}(H,G)$, et c'est un \deleted{\unsure{schéma}}
$S$-préschéma constant tordu. En particulier, $\underline{\operatorname{Aut}}_{S\text{-gr}}(H)$ est représentable
et est un groupe constant tordu (en général non commutatif) sur $S$.

Cela résulte de 5.8. et de Exp VIII\deleted{\unsure{.}} 1.6.

\uline{Proposition} 5.11. Soient $S$ un préschéma, $R$ un groupe constant tordu
sur $S$, $H=D_S(R)$ le groupe de type multiplicatif qu'il définit. Considérons les conditions suivantes:

(i) $H$ est isotrivial, (i.e. $R$ est isotrivial).

(ii) $R$ est \deleted{somme disjointe de}\inserted{réunion de} sous-préschémas\inserted{à la fois ouverts et fermés} $R_i$, qui sont quasi-compacts
(et alors nécessairement finis) sur $S$.

(iii) Les composantes connexes de $R$ sont finies sur $S$.

Alors $(i)\Rightarrow(ii)\Rightarrow(iii)$, on a $\inserted{(i)\Leftrightarrow}(ii)\Leftrightarrow(iii)$ si $S$ est localement
noethérien, enfin $(i)\Leftrightarrow(ii)$ si $R$ est \deleted{\unsure{de}} à engendrement fini i.e.
si $H$ est de type fini, \inserted{du moins si $S$ est quasi-compact ou si ses
composantes connexes sont ouvertes.}

\deleted{Démonstration. $(i)\unsure{-|\Rightarrow}(ii)$ Soit $S'-S$ fini étale surjectif qui
splitte $H$ donc $R$, de sorte que $R'$ est \unsure{somm}}

Décomposant d'abord $S$ en préschéma somme de préschémas $S_i$ sur chacun desquels $H$ est de type constant (cf Exp IX\deleted{\unsure{.}} 1), on
est ramené au cas où $H$ \deleted{est de} donc $R$ est de type constant. Nous aurons
besoin\deleted{\unsure{(pour (ii))}} d'un

\deleted{\uline{Lemme} 5.11.}
}
\mspage{17}

